% Compile with LuaLaTeX
\documentclass[11pt,a4paper]{article}


% ==========================================
% 1. MODERN FONTS & LANGUAGES
% ==========================================
\usepackage{fontspec}
% Tip: For Romance languages and IPA, "Charis SIL" or "Doulos SIL" are the best fonts. 
% If you have them installed, uncomment the next line:
% \setmainfont{Charis SIL} 

\usepackage{polyglossia}
\setmainlanguage{english} % change to german if needed
% Load the Romance languages you study
\setotherlanguages{french,spanish,italian,portuguese,romanian,latin}

% ==========================================
% 1b. TYPOGRAPHY & QUOTES
% ==========================================
\usepackage{microtype}   % Better line breaks, especially helpful for multilingual docs
\usepackage{csquotes}    % Language-aware quotes: \enquote{...} picks the right marks
\usepackage{booktabs}
\usepackage{graphicx}

% ==========================================
% 2. LINGUISTICS & PHONETICS
% ==========================================
\usepackage{tipauni}     % Modern, LuaLaTeX-safe IPA symbols
\usepackage{gb4e}        % For interlinear glosses (standard in linguistics)
\noautomath              % CRITICAL: Prevents gb4e from crashing with modern packages
\usepackage{forest}      % For drawing modern syntax trees

% ==========================================
% 3. MATH & LAYOUT
% ==========================================
\usepackage{amsmath,amssymb,mathtools}
\usepackage{geometry}
\geometry{margin=2.5cm}
\usepackage{hyperref}
\usepackage{xcolor}
\usepackage{siunitx}

\usepackage{enumitem}


\title{Computational Linguistics \\ LaTeX Cheat-Sheet}
\author{}
\date{}

\begin{document}

\maketitle
\tableofcontents
\vspace{1em}

%-------------------------------------------------
\section{Basic Math \& Statistics}

\begin{itemize}[leftmargin=*]
  \item Inline math: \texttt{\$P(w)\$} $\rightarrow$ $P(w)$
  \item Display math: \texttt{\textbackslash[ ... \textbackslash]} or \texttt{equation} environment
\end{itemize}

\subsection{Most common symbols}

\begin{center}
\begin{tabular}{ll@{\hspace{1.8cm}}ll}
\toprule
\textbf{Command} & \textbf{Result} & \textbf{Command} & \textbf{Result} \\
\midrule
\texttt{\textbackslash alpha}   & $\alpha$   & \texttt{\textbackslash beta}    & $\beta$    \\
\texttt{\textbackslash lambda}  & $\lambda$  & \texttt{\textbackslash theta}   & $\theta$   \\
\texttt{\textbackslash mu}      & $\mu$      & \texttt{\textbackslash sigma}   & $\sigma$   \\
\texttt{\textbackslash sum}     & $\sum$     & \texttt{\textbackslash prod}    & $\prod$    \\
\texttt{\textbackslash frac\{a\}\{b\}} & $\frac{a}{b}$ & \texttt{\textbackslash sqrt\{x\}} & $\sqrt{x}$ \\
\texttt{\textbackslash log}     & $\log$     & \texttt{\textbackslash exp}     & $\mathrm{exp}$ \\
\texttt{\textbackslash infty}   & $\infty$   & \texttt{\textbackslash approx}  & $\approx$  \\
\texttt{\textbackslash leq}     & $\leq$     & \texttt{\textbackslash geq}     & $\geq$     \\
\texttt{\textbackslash times}   & $\times$   & \texttt{\textbackslash cdot}    & $\cdot$    \\
\texttt{\textbackslash rightarrow} & $\rightarrow$ & \texttt{\textbackslash Rightarrow} & $\Rightarrow$ \\
\bottomrule
\end{tabular}
\end{center}

\subsection{Useful statistical examples}

\begin{verbatim}
P(w_i | w_{i-1}) = \frac{C(w_{i-1}, w_i)}{C(w_{i-1})}

\text{Precision} = \frac{TP}{TP + FP}
\text{Recall}    = \frac{TP}{TP + FN}
\text{F1}        = 2 \cdot \frac{P \cdot R}{P + R}
\end{verbatim}

%-------------------------------------------------
\section{IPA Phonetic Symbols (\texttt{tipauni})}

\begin{itemize}[leftmargin=*]
  \item Preferred modern way: type the Unicode character directly (if your font supports it).
  \item Classic shortcuts still work with \texttt{\textbackslash textipa\{...\}}.
\end{itemize}

\begin{center}
\begin{tabular}{ll@{\hspace{1.2cm}}ll}
\toprule
\textbf{Command / Unicode} & \textbf{Result} & \textbf{Command / Unicode} & \textbf{Result} \\
\midrule
\texttt{\textbackslash textipa\{[aI]\}} / aɪ     & \textipa{[aɪ]}  & \texttt{\textbackslash textipa\{[E]\}} / ɛ      & \textipa{[ɛ]} \\
\texttt{\textbackslash textipa\{[S]\}} / ʃ       & \textipa{[ʃ]}   & \texttt{\textbackslash textipa\{[Z]\}} / ʒ      & \textipa{[ʒ]} \\
\texttt{\textbackslash textipa\{[T]\}} / θ       & \textipa{[θ]}   & \texttt{\textbackslash textipa\{[D]\}} / ð      & \textipa{[ð]} \\
\texttt{\textbackslash textipa\{[N]\}} / ŋ       & \textipa{[ŋ]}   & \texttt{\textbackslash textipa\{[@]\}} / ə      & \textipa{[ə]} \\
\texttt{\textbackslash textipa\{[p\~{}]\}} / p̃   & \textipa{[p̃]}  & \texttt{\textbackslash textipa\{[e:]\}} / eː    & \textipa{[eː]} \\
\bottomrule
\end{tabular}
\end{center}

%-------------------------------------------------
\section{Linguistic Examples \& Glosses (\texttt{gb4e})}

\begin{verbatim}
\begin{exe}
  \ex \gll  Elle a mangé une pomme. \\
            She has eaten an  apple. \\
      \glt  `She ate an apple.'
\end{exe}
\end{verbatim}

Result:

\begin{exe}
  \ex \gll  Elle a mangé une pomme. \\
            She has eaten an  apple. \\
      \glt  `She ate an apple.'
\end{exe}

\vspace{0.8em}
More examples in one block:

\begin{verbatim}
\begin{exe}
  \ex \gll ... \\
            ... \\
      \glt  `...'
  \ex \gll ... \\
            ... \\
      \glt  `...'
\end{exe}
\end{verbatim}

%-------------------------------------------------
\section{Syntax Trees (\texttt{forest})}

\begin{verbatim}
\begin{forest}
  [S
    [NP [Det the] [N cat]]
    [VP [V sleeps]]
  ]
\end{forest}
\end{verbatim}

%-------------------------------------------------
\section{Multilingual Text (\texttt{polyglossia} + \texttt{csquotes})}

\begin{verbatim}
\textfrench{Bonjour !}
\textspanish{¡Hola!}
\textitalian{Ciao!}
\textportuguese{Olá!}

\enquote{This automatically uses the correct quotation marks.}
\end{verbatim}

%-------------------------------------------------
\section{Extension Packages}

Load these as you grow into them --- all are well-maintained and ship with \texttt{texlive-full}.

\begin{itemize}[leftmargin=*]
  \item \texttt{microtype} --- better line breaks, especially in multilingual documents.
  \item \texttt{csquotes} --- language-aware quotes via \verb|\enquote{...}|.
  \item \texttt{siunitx} --- aligned numeric tables and in-text numbers (\verb|\num{1.23e-4}|, \verb|\SI{95}{\percent}|).
  \item \texttt{pgfplots} --- vector-based plots (including Zipf-style log-log plots).
\end{itemize}

%-------------------------------------------------
\section{Quick Tips}

\begin{itemize}[leftmargin=*]
  \item Always compile with \textbf{LuaLaTeX}.
  \item Use \texttt{latexmk -lualatex -pdf file.tex} to handle multi-pass compilation + Biber automatically.
  \item For bibliography in linguistics: \texttt{biblatex} with \texttt{style=unified} (LSA) or \texttt{style=authoryear}.
  \item Need a symbol? Use Detexify or the Comprehensive LaTeX Symbol List.
\end{itemize}

\end{document}